% ============================================================
%  BAB 5 – PEMBAHASAN
%  Compile standalone:  pdflatex bab-5.tex
%  Compile as part of thesis: pdflatex main.tex
% ============================================================
\documentclass[../../thesis]{subfiles}
\usepackage{hyphenat}
\usepackage{iftex}

% ── Encoding, Language & Fonts ───────────────────────────────
\ifXeTeX
    \usepackage{polyglossia}
    \setmainlanguage{bahasai}
    \usepackage{fontspec}
    \setmainfont{TeX Gyre Termes}
    \setsansfont{TeX Gyre Heros}
    \setmonofont[Scale=0.88]{TeX Gyre Cursor}
\else
    \usepackage[utf8]{inputenc}
    \usepackage[T1]{fontenc}
    \usepackage[indonesian]{babel}
    \usepackage{mathptmx}
    \usepackage{helvet}
    \usepackage{courier}
\fi

% ── Page geometry (A4, UI thesis margins) ────────────────────
\usepackage[
  a4paper,
  top    = 2.54cm,
  bottom = 2.54cm,
  left   = 3.17cm,
  right  = 2.54cm
]{geometry}

% ── Line spacing & paragraph ─────────────────────────────────
\usepackage{setspace}
\onehalfspacing
\setlength{\parindent}{1.27cm}
\setlength{\parskip}{0pt}

% ── Microtype & justification ────────────────────────────────
\usepackage{microtype}
\sloppy

% ── Headings ─────────────────────────────────────────────────
\usepackage{titlesec}

\titleformat{\chapter}[block]
  {\centering\rmfamily\bfseries\large}{}
  {0pt}{\MakeUppercase}
\titlespacing*{\chapter}{0pt}{0pt}{1.5em}

\titleformat{\section}
  {\rmfamily\bfseries\normalsize}{\thesection}{1em}{}
\titlespacing*{\section}{0pt}{1.2em}{0.6em}

\titleformat{\subsection}
  {\rmfamily\bfseries\normalsize}{\thesubsection}{1em}{}
\titlespacing*{\subsection}{0pt}{1em}{0.5em}

\titleformat{\subsubsection}
  {\rmfamily\bfseries\normalsize}{\thesubsubsection}{1em}{}
\titlespacing*{\subsubsection}{0pt}{0.8em}{0.4em}

% ── Tables ───────────────────────────────────────────────────
\usepackage{longtable, booktabs, array, multirow}
\usepackage{calc}
\usepackage{tabularx}

% ── Math ─────────────────────────────────────────────────────
\usepackage{amsmath, amssymb}

% ── Graphics ─────────────────────────────────────────────────
\usepackage{graphicx}
\makeatletter
\def\maxwidth{\ifdim\Gin@nat@width>\linewidth\linewidth\else\Gin@nat@width\fi}
\makeatother
\setkeys{Gin}{width=\maxwidth,keepaspectratio}

% ── TikZ + pgfplots (native bridge graph, sweep & relation charts) ──
\usepackage{tikz}
\usetikzlibrary{calc, positioning, fit, backgrounds}
\usepackage{pgfplots}
\usepgfplotslibrary{groupplots}
\pgfplotsset{compat=1.18}

% ── Lists ────────────────────────────────────────────────────
\usepackage{enumitem}
\setlist[enumerate,1]{
  label      = \arabic*.,
  leftmargin = 1.27cm,
  itemsep    = 4pt,
  parsep     = 0pt,
  topsep     = 4pt
}
\setlist[enumerate,2]{
  label      = \alph*.,
  leftmargin = 1.27cm,
  itemsep    = 3pt,
  parsep     = 0pt,
  topsep     = 2pt
}

% ── Hyperlinks ───────────────────────────────────────────────
\usepackage[hidelinks]{hyperref}

% Bibliography setup is inherited from main.tex (biblatex + references.bib)
% when this file is compiled standalone via the subfiles class.

% ============================================================
\begin{document}

% When standalone: set counter so \chapter{} renders as "BAB 5".
% When in main.tex: main.tex sets the counter before \subfile{bab-5}.
\ifSubfilesClassLoaded{}{\setcounter{chapter}{4}}

\chapter{Pembahasan}\label{sec:hasil-evaluasi}

Setelah artefak selesai dibangun, pertanyaannya kini adalah apakah ia
benar-benar menjembatani fragmentasi yang melatari penelitian ini. Sebagai tahap evaluasi
dalam kerangka \textit{Design Science Research} (DSR), bab ini menghadapkan
\textit{Knowledge Graph} (KG) yang dihasilkan pada pengujian untuk menilai
sejauh mana relasi yang ditemukan sahih secara struktural sekaligus bermakna
secara pedagogis. Pembahasan bergerak dari bukti topologis perubahan graf
menuju penilaian pakar atas ketepatan relasi, menelaah pola kesalahan model
LLM yang terungkap di baliknya, dan berakhir pada perbandingan karakteristik
ketiga disiplin sains.

% ────────────────────────────────────────────────────────────
\section{Hasil Analisis Struktural Graf}\label{subsec:struktural}
% ────────────────────────────────────────────────────────────

Analisis struktural bertujuan untuk mengevaluasi kualitas interkoneksi
jaringan pengetahuan secara matematis, khususnya untuk melihat
perbandingan struktur graf sebelum \textit{Knowledge Graph Completion} (pre-KGC)
dan sesudahnya (post-KGC). Seluruh perhitungan pada bagian ini
dieksekusi melalui kueri \textit{Cypher} di Neo4j dengan berpedoman pada
formulasi metrik topologi yang telah dijelaskan pada Subbab~\ref{subsec:metrik-topologi}.
Perbandingan topologi graf sebelum dan sesudah KGC ditunjukkan pada
Gambar~\ref{fig:kg-topologi}: pada kondisi pra-KGC tampak tiga graf
yang terpisah, mencerminkan ketiga komunitas mata pelajaran yang
terisolasi, sedangkan pada kondisi pasca-KGC ketiga graf tersebut saling
terhubung oleh relasi lintas-buku.

\begin{figure}[htbp]
  \centering
  \includegraphics[width=0.82\linewidth]{assets/pics/kg/pre-kgc.png}\\[0.2em]
  {\small (a) Pra-KGC}\\[0.8em]
  \includegraphics[width=0.82\linewidth]{assets/pics/kg/post-kgc.png}\\[0.2em]
  {\small (b) Pasca-KGC}
  \caption{Perbandingan topologi \textit{knowledge graph} (a) sebelum dan
    (b) sesudah KGC; relasi lintas-buku menjembatani ketiga komunitas mapel.}
  \label{fig:kg-topologi}
\end{figure}

\begin{table}[htbp]
\centering
\caption{Perbandingan metrik struktural pre-KGC dan post-KGC.}
\label{tab:struktural}
\begin{tabular}{@{}
  >{\raggedright\arraybackslash}p{0.38\linewidth}
  >{\centering\arraybackslash}p{0.27\linewidth}
  >{\centering\arraybackslash}p{0.27\linewidth}@{}}
\toprule
\textbf{Metrik Struktural} & \textbf{Pre-KGC} & \textbf{Post-KGC} \\
\midrule
AD                & 0{,}7384 & 1{,}4651 \\
\textit{Graph Density} & 0{,}00108 & 0{,}00214 \\
\textit{Modularity}               & 0{,}6659 & 0{,}1580 \\
\bottomrule
\end{tabular}
\end{table}

\begin{table}[htbp]
\centering
\caption{Uji ketahanan metrik struktural dengan dan tanpa \textit{edge}
berlabel \texttt{wrong}.}
\label{tab:struktural-robustness}
\begin{tabular}{@{}
  >{\raggedright\arraybackslash}p{0.30\linewidth}
  >{\centering\arraybackslash}p{0.16\linewidth}
  >{\centering\arraybackslash}p{0.16\linewidth}
  >{\centering\arraybackslash}p{0.16\linewidth}
  >{\centering\arraybackslash}p{0.16\linewidth}@{}}
\toprule
& \multicolumn{2}{c}{\textbf{\textit{Wrong} Dipertahankan}} &
  \multicolumn{2}{c}{\textbf{\textit{Wrong} Dikeluarkan}} \\
\cmidrule(lr){2-3}\cmidrule(l){4-5}
\textbf{Metrik} & \textbf{Pre-KGC} & \textbf{Post-KGC} &
\textbf{Pre-KGC} & \textbf{Post-KGC} \\
\midrule
ADC & 0{,}7384 & 1{,}4651 & 0{,}7093 & 1{,}4360 \\
\textit{Graph Density} & 0{,}00108 & 0{,}00214 & 0{,}00103 & 0{,}00209 \\
\textit{Modularity} & 0{,}6659 & 0{,}1580 & 0{,}6654 & 0{,}1489 \\
\textit{Concept}$\rightarrow$\textit{Concept} \textit{edge}
& 127 & 252 & 122 & 247 \\
\bottomrule
\end{tabular}
\end{table}

\begin{figure}[htbp]
\centering
\begin{tikzpicture}
\begin{groupplot}[
  group style={
    group size=3 by 1,
    horizontal sep=1.35cm
  },
  width=0.32\linewidth,
  height=5.2cm,
  ybar,
  /pgf/bar width=8pt,
  ymin=0,
  enlarge x limits=0.38,
  symbolic x coords={Pre-KGC,Post-KGC},
  xtick=data,
  xticklabel style={font=\scriptsize},
  yticklabel style={font=\scriptsize},
  label style={font=\scriptsize},
  tick label style={font=\scriptsize},
  title style={font=\small\bfseries},
  nodes near coords,
  nodes near coords align={vertical},
  every node near coord/.append style={
    font=\tiny,
    yshift=2pt
  }
]

\nextgroupplot[
  title={AD},
  ylabel={Nilai},
  ymax=1.7,
  nodes near coords={\pgfmathprintnumber[fixed,precision=2]{\pgfplotspointmeta}}
]
\addplot+[fill=blue!25, draw=blue] coordinates {(Pre-KGC,0.7384) (Post-KGC,1.4651)};
\addplot+[fill=red!25, draw=red] coordinates {(Pre-KGC,0.7093) (Post-KGC,1.4360)};

\nextgroupplot[
  title={\textit{Graph Density}},
  ymax=0.00245,
  scaled y ticks=false,
  y tick label style={
    /pgf/number format/fixed,
    /pgf/number format/precision=4,
    font=\scriptsize
  },
  nodes near coords={\pgfmathprintnumber[sci,precision=2]{\pgfplotspointmeta}},
  legend style={
    font=\scriptsize,
    at={(0.5,-0.26)},
    anchor=north,
    legend columns=2,
    draw=none
  }
]
\addplot+[fill=blue!25, draw=blue] coordinates {(Pre-KGC,0.00108) (Post-KGC,0.00214)};
\addplot+[fill=red!25, draw=red] coordinates {(Pre-KGC,0.00103) (Post-KGC,0.00209)};
\legend{\textit{Wrong} dipertahankan,\textit{Wrong} dikeluarkan}

\nextgroupplot[
  title={\textit{Modularity}},
  ymax=0.76,
  nodes near coords={\pgfmathprintnumber[fixed,precision=2]{\pgfplotspointmeta}}
]
\addplot+[fill=blue!25, draw=blue] coordinates {(Pre-KGC,0.6659) (Post-KGC,0.1580)};
\addplot+[fill=red!25, draw=red] coordinates {(Pre-KGC,0.6654) (Post-KGC,0.1489)};

\end{groupplot}
\end{tikzpicture}
\caption{Perbandingan metrik struktural ketika relasi berlabel \texttt{wrong}
dipertahankan dan dikeluarkan dari subgraf \texttt{Concept}$\rightarrow$\texttt{Concept}.}
\label{fig:robustness-wrong-comparison}
\end{figure}

\begin{enumerate}
  \item \textbf{Evaluasi \textit{Average Degree} (AD) dan Kepadatan.}

  Berdasarkan Tabel~\ref{tab:struktural}, AD antar-konsep hampir berlipat ganda
  dari 0{,}74 (pre-KGC) menjadi 1{,}47 (post-KGC), naik 98{,}4\%. Peningkatan
  tajam ini mencerminkan bahwa \textit{pipeline} KGC berhasil menambahkan 125
  relasi lintas-buku baru yang secara langsung menghubungkan \textit{node}
  \texttt{Concept} antar-mapel, meningkatkan kepadatan interkoneksi semantik
  yang sebelumnya tidak ada. Perlu dicatat bahwa seluruh metrik struktural pada
  bagian ini dihitung atas \textbf{125 \textit{edge}} keluaran \textit{pipeline}
  pada konfigurasi final---termasuk 8 \textit{edge} berkeyakinan rendah
  (\textit{confidence} $< 0{,}5$) yang tidak diikutkan dalam survei pakar
  (Subbab~\ref{subsec:kualitatif})---sehingga himpunan ini berbeda dari
  \textbf{117 \textit{edge}} yang menjadi basis validasi pakar pada
  Subbab~\ref{subsec:validasi-pakar}.
  Kueri Cypher untuk menghitung metrik struktural ini diberikan pada
  \ref{lamp:cypher-metrik}, sedangkan daftar lengkap \textit{edge}
  lintas-buku yang dievaluasi pada Fase~2 beserta hasil validasinya diberikan
  pada \ref{lamp:lintas-buku}.
  \textit{Graph Density} antar-konsep, yakni kepadatan relasi pada subgraf
  \texttt{Concept} (lihat Subbab~\ref{subsec:metrik-topologi}), juga meningkat
  secara proporsional dari 0{,}00108 menjadi 0{,}00214 (hampir dua kali lipat),
  konsisten dengan penambahan \textit{edge} lintas-buku tersebut. Nilai densitas
  yang masih sangat rendah pada kedua fase merupakan karakteristik umum
  \textit{knowledge graph} berukuran sedang-besar, karena hanya sebagian kecil
  dari seluruh kemungkinan pasangan konsep yang memiliki relasi bermakna
  \parencite{hogan_knowledge_2022}.
  Untuk memastikan bahwa pola ini tidak bergantung pada relasi yang dinilai salah
  oleh pakar, Tabel~\ref{tab:struktural-robustness} menyajikan perhitungan ulang
  setelah seluruh \textit{edge} berlabel \texttt{wrong} dikeluarkan dari subgraf
  \texttt{Concept}$\rightarrow$\texttt{Concept}. Hasilnya tetap searah: AD
  meningkat dari 0{,}7093 menjadi 1{,}4360, sedangkan densitas meningkat dari
  0{,}00103 menjadi 0{,}00209. Visualisasi pada
  Gambar~\ref{fig:robustness-wrong-comparison} memperlihatkan bahwa selisih
  antara kondisi \textit{wrong} dipertahankan dan dikeluarkan relatif kecil,
  sementara arah perubahan pre-KGC ke post-KGC tetap konsisten. Dengan demikian,
  kesimpulan bahwa KGC meningkatkan kepadatan hubungan antar-konsep tidak
  bergantung pada keberadaan \textit{edge} salah; daftar relasi berlabel
  \texttt{wrong} disajikan pada \ref{lamp:wrong-triples}.

  \item \textbf{Evaluasi \textit{Modularity}.}

  \textit{Modularity} di sini dihitung atas partisi komunitas \textbf{tetap}
  berdasarkan mata pelajaran (Fisika, Kimia, Biologi), sesuai dengan
  formulasi pada Subbab~\ref{subsec:metrik-topologi}. Dengan demikian, perubahan
  nilai $Q$ murni mencerminkan bertambahnya \textit{edge} lintas-mapel.
  Analisis struktur komunitas menggunakan metrik \textit{Modularity} pada fase
  pre-KGC mencatatkan nilai 0{,}666, yang mengindikasikan struktur komunitas yang
  sangat kuat --- ketiga komunitas mapel (Fisika, Kimia, Biologi) hampir sepenuhnya
  terisolasi satu sama lain, sejalan dengan kondisi \textit{siloed learning} yang
  diidentifikasi sebagai masalah penelitian. Setelah KGC, nilai \textit{Modularity}
  turun secara signifikan menjadi 0{,}158, menandakan bahwa relasi lintas-buku yang
  ditambahkan berhasil menjembatani ketiga komunitas tersebut dan mengurangi
  keterpisahan topologis antar-disiplin. Penurunan ini merupakan indikator
  keberhasilan utama \textit{pipeline}, karena penurunan \textit{Modularity}
  mencerminkan peningkatan interkoneksi lintas-komunitas sebagaimana dirumuskan
  pada Subbab~\ref{subsec:metrik-topologi}. Setelah \textit{edge} berlabel
  \texttt{wrong} dikeluarkan, \textit{Modularity} tetap turun dari 0{,}6654
  menjadi 0{,}1489, sehingga arah kesimpulan struktural tetap sama.

  \item \textbf{Evaluasi Metrik Ekstraksi dan Hubungan.}

  Di luar evaluasi topologi di atas, kualitas hasil ekstraksi LLM secara
  spesifik diukur menggunakan metrik \textit{Description Completeness} (DC),
  \textit{Empty SubKonsep Rate} (ESR), dan
  \textit{Typed Relation Ratio} (TRR) sebagaimana didefinisikan pada
  Subbab~\ref{subsec:metrik-ekstraksi} dan~\ref{subsec:metrik-hubungan}.
  Sistem mencatatkan nilai DC sebesar 93{,}07\% --- identik pada
  fase pre-KGC maupun post-KGC karena DC mengukur kualitas deskripsi \textit{node} yang
  tidak berubah akibat penambahan \textit{edge} KGC. Metrik ESR tercatat 0\%
  (0,00) pada kedua fase --- jauh di bawah ambang maksimal 10\% --- yang
  menandakan bahwa seluruh entitas induk (\texttt{:Subtopic}) berhasil
  didekomposisi menjadi setidaknya satu konsep turunan, sehingga tidak ada
  ``\textit{node} buntu'' dalam graf. Nilai TRR meningkat dari
  69{,}69\% (pre-KGC) menjadi 74{,}86\% (post-KGC), mencerminkan bahwa KGC
  menambahkan \textit{edge} bertipe semantik spesifik sehingga proporsi relasi
  bertipe meningkat relatif terhadap total \textit{edge} dalam graf.
\end{enumerate}

% ────────────────────────────────────────────────────────────
\section{Hasil Validasi Pakar}\label{subsec:validasi-pakar}
% ────────────────────────────────────────────────────────────

Performa sistem dalam mengekstrak dan memprediksi relasi dinilai langsung oleh
panel enam pakar menggunakan instrumen KG Review App. Sesuai desain studi dua
fase, hasil penilaian dipisahkan menjadi validasi intra-mapel (Fase 1) dan
inter-mapel (Fase 2). Pada Fase~1, setiap mata pelajaran dinilai oleh dua
pakar independen yang dilabeli ``Expert~1'' dan ``Expert~2'' pada
Tabel~\ref{tab:validasi-mapel}; label ini bersifat lokal per mata
pelajaran, sehingga ``Expert~1 Biologi'' bukan orang yang sama dengan
``Expert~1 Fisika'' maupun ``Expert~1 Kimia''---total enam pakar berbeda
terlibat. Data dianalisis menggunakan metrik \textit{Fuzzy
Precision} dan \textit{Fuzzy Recall} \parencite{harju_evaluating_2023}, F1-Score,
serta uji kesepakatan Gwet's AC1, sebagaimana didefinisikan pada
Subbab~\ref{subsec:validasi-manusia}.

{\setlength{\tabcolsep}{4pt}
\begin{table}[htbp]
\centering
\caption{Hasil validasi Fase~1 (intra-mapel) per pakar dan konsensus final.
  P/R = \textit{Fuzzy Precision}/\textit{Recall} (bobot pada
  Subbab~\ref{subsec:validasi-manusia}); semua \textit{triple} yang ditinjau,
  termasuk berlabel Salah, masuk ke penyebut presisi. Kolom \textbf{Hilang}
  adalah jumlah \textit{missing triples}, yakni relasi esensial yang dilaporkan
  pakar sebagai terlewat oleh sistem; relasi ini tidak termasuk dalam kolom
  Ditinjau dan menjadi \textit{false negative} yang menurunkan \textit{Recall}
  (rincian pada \ref{lamp:missing-triples}).}
\label{tab:validasi-mapel}
\begin{tabular}{@{}l l c c c c c c c c@{}}
\toprule
\textbf{Mapel} & \textbf{Penilai} & \textbf{Ditinjau} & \textbf{Benar} & \textbf{Parsial} & \textbf{KK} & \textbf{Salah} & \textbf{Hilang} & \textbf{P} & \textbf{R} \\
\midrule
Biologi  & Expert 1      & 151 & 140 & 6  & 3  & 2  & 0  & 0{,}952 & 0{,}965 \\
         & Expert 2      & 151 & 143 & 6  & 2  & 0  & 0  & 0{,}970 & 0{,}970 \\
         & \textit{Final}& 151 & 144 & 4  & 2  & 1  & 0  & \textbf{0{,}970} & \textbf{0{,}977} \\
\midrule
Fisika   & Expert 1      & 240 & 239 & 0  & 1  & 0  & 15 & 0{,}997 & 0{,}938 \\
         & Expert 2      & 240 & 186 & 40 & 11 & 3  & 0  & 0{,}870 & 0{,}881 \\
         & \textit{Final}& 240 & 219 & 19 & 0  & 2  & 10 & \textbf{0{,}952} & \textbf{0{,}921} \\
\midrule
Kimia    & Expert 1      & 166 & 132 & 7  & 2  & 25 & 0  & 0{,}819 & 0{,}965 \\
         & Expert 2      & 166 & 136 & 6  & 19 & 5  & 29 & 0{,}866 & 0{,}757 \\
         & \textit{Final}& 166 & 140 & 6  & 2  & 18 & 19 & \textbf{0{,}864} & \textbf{0{,}859} \\
\midrule
Keseluruhan & Expert 1   & 557 & 511 & 13 & 6  & 27 & 15 & 0{,}932 & 0{,}952 \\
            & Expert 2   & 557 & 465 & 52 & 32 & 8  & 29 & 0{,}896 & 0{,}863 \\
            & \textit{Final} & 557 & 503 & 29 & 4 & 21 & 29 & \textbf{0{,}931} & \textbf{0{,}918} \\
\bottomrule
\end{tabular}
\end{table}
}

\begin{table}[htbp]
\centering
\caption{Uji kesepakatan antarpakar Gwet's AC1 per mata pelajaran (Fase 1).}
\label{tab:ac1}
\begin{tabular}{@{}
  >{\raggedright\arraybackslash}p{0.18\linewidth}
  >{\centering\arraybackslash}p{0.14\linewidth}
  >{\centering\arraybackslash}p{0.14\linewidth}
  >{\centering\arraybackslash}p{0.14\linewidth}
  >{\centering\arraybackslash}p{0.14\linewidth}
  >{\centering\arraybackslash}p{0.14\linewidth}@{}}
\toprule
\textbf{Mapel} & \textbf{Item} & \textbf{Obs. Agreement} & \textbf{Chance Agr.} & \textbf{Gwet's AC1} & \textbf{Interpretasi} \\
\midrule
Biologi & 151 & 0{,}934 & 0{,}040 & \textbf{0{,}931} & Sangat baik \\
Fisika  & 240 & 0{,}771 & 0{,}069 & \textbf{0{,}754} & Baik \\
Kimia   & 166 & 0{,}651 & 0{,}112 & \textbf{0{,}607} & Baik \\
\bottomrule
\end{tabular}
\end{table}

\begin{table}[htbp]
\centering
\caption{Ringkasan metrik per fase evaluasi.}
\label{tab:validasi-metrik}
\begin{tabular}{@{}
  >{\raggedright\arraybackslash}p{0.20\linewidth}
  >{\centering\arraybackslash}p{0.18\linewidth}
  >{\centering\arraybackslash}p{0.18\linewidth}
  >{\centering\arraybackslash}p{0.18\linewidth}
  >{\centering\arraybackslash}p{0.18\linewidth}@{}}
\toprule
\textbf{Fase} & \textbf{Fuzzy Precision / Validitas} & \textbf{Fuzzy Recall} & \textbf{F1-Score} & \textbf{Gwet's AC1} \\
\midrule
Fase 1 (konsensus) & 0{,}931 & 0{,}918 & 0{,}924 & lihat Tabel~\ref{tab:ac1} \\
Fase 2 (panel diskusi) & 0{,}991 & — & — & — \\
\bottomrule
\end{tabular}
\vspace{0.35em}
\begin{minipage}{0.88\linewidth}
\footnotesize
\textit{Catatan.} Pada Fase~1, kolom pertama adalah \textit{Fuzzy Precision}.
Pada Fase~2, kolom yang sama dibaca sebagai validitas relasi terprediksi, yaitu
proporsi kandidat relasi lintas-mapel yang dinilai valid oleh panel.
\textit{Fuzzy Recall}, F1-Score, dan Gwet's AC1 tidak dihitung
karena evaluasi dilakukan melalui panel diskusi dan tidak menyediakan penyebut
relasi benar yang menyeluruh (lihat Subbab~\ref{subsec:metrik-fase2}).
\end{minipage}
\end{table}

\begin{enumerate}
  \item \textbf{Kinerja Ekstraksi Intra-mapel (Fuzzy Precision, Fuzzy Recall, F1-Score).}

  Ketiga metrik ini tidak diperlakukan sebagai akurasi klasifikasi sederhana,
  melainkan sebagai ukuran kinerja ekstraksi relasi: \textit{Fuzzy Precision}
  menunjukkan proporsi relasi hasil sistem yang dinilai benar atau sebagian
  benar oleh pakar, \textit{Fuzzy Recall} menunjukkan sejauh mana sistem tidak
  melewatkan relasi esensial yang kemudian diusulkan sebagai
  \textit{missing triples}, dan F1-Score merangkum keseimbangan keduanya.
  Tabel~\ref{tab:validasi-metrik}, yang berjudul Ringkasan metrik per fase evaluasi,
  menunjukkan bahwa Fase~1 menghasilkan \textit{Fuzzy Precision} 93{,}1\%,
  \textit{Fuzzy Recall} 91{,}8\%, dan F1-Score 92{,}4\%. Rincian per mata
  pelajaran pada Tabel~\ref{tab:validasi-mapel}, yang berjudul Hasil validasi
  Fase~1 per pakar dan konsensus final, menunjukkan bahwa
  kemampuan LLM dalam mengekstraksi relasi intra-mapel tergolong baik. Biologi
  mencatatkan F1 tertinggi (0{,}973), diikuti Fisika (0{,}936), dan Kimia (0{,}862).

  Pada level keseluruhan, \textit{Fuzzy Precision} melampaui \textit{Fuzzy Recall} sebesar
  1{,}3 poin. Selisih ini melebar pada konsensus final Fisika (3{,}1 poin) dan
  terutama pada penilaian Expert 2 Kimia (10{,}9 poin). Penyebab utamanya adalah
  \textit{missing triples} (10 relasi pada konsensus final Fisika dan 29 relasi
  yang dilaporkan Expert 2 Kimia) yang menandakan sistem masih luput
  mengekstrak sebagian relasi kuantitatif dan terminologi proses yang esensial.
  Rincian seluruh \textit{missing triples} yang diusulkan pakar beserta putusan
  penengahan \textit{LLM-as-a-judge} per mata pelajaran diberikan pada
  \ref{lamp:missing-triples}.
  Perbedaan standar antarpakar juga tampak jelas pada Fisika (Expert 1 vs.\
  Expert 2: P 0{,}997 vs.\ 0{,}870; R 0{,}938 vs.\ 0{,}881), dan dikonfirmasi
  oleh nilai AC1 Fisika (0{,}754) yang lebih rendah daripada Biologi (0{,}931).

  Pada Fase 2 (inter-mapel), validasi panel diskusi atas 117 relasi lintas-buku
  menghasilkan validitas relasi terprediksi 99{,}1\% (116/117 relasi
  dikonfirmasi valid; satu relasi ditolak). Nilai ini berbeda dari
  \textit{Fuzzy Precision} Fase~1 karena Fase~2 memakai keputusan panel tunggal
  dan tidak menggunakan skala parsial. Metrik \textit{Fuzzy Recall} dan F1 tidak dapat dihitung
  karena tidak tersedia informasi jumlah total relasi lintas-buku yang benar
  (penyebut untuk menghitung cakupan relasi benar) dalam konfigurasi panel tunggal ini. Karena validasi
  dilakukan melalui diskusi panel (bukan dua pakar independen), uji Gwet's AC1
  tidak berlaku.

  \item \textbf{Uji Kesepakatan Pakar (Inter-rater Reliability).}

  Sesuai dengan pilihan metodologi yang ditetapkan pada Bab 3 (dan latar belakang
  teoritis Gwet's AC1 pada Subbab~\ref{subsec:validasi-manusia}), uji kesepakatan antarpakar menggunakan
  Gwet's AC1 \parencite{gwet_handbook_2014}. Tabel~\ref{tab:ac1}, yang
  berjudul Uji kesepakatan antarpakar Gwet's AC1 per mata pelajaran,
  menunjukkan variasi tingkat kesepakatan
  yang signifikan antar mata pelajaran. Biologi mencatatkan AC1 tertinggi dan
  tergolong sangat baik (0{,}931), mengindikasikan kesepakatan yang hampir
  sempurna antara kedua pakar. Hal ini konsisten dengan sifat konten Biologi
  yang lebih deskriptif dan definitif. Fisika (0{,}754) dan Kimia (0{,}607)
  sama-sama berada pada kategori baik, dengan Kimia sebagai yang terendah.
  Nilai AC1 Kimia yang paling rendah selaras dengan pola kesalahan pada
  Subbab~\ref{subsubsec:error-analysis}: terminologi proses kimia (seperti elektrorefining vs.\
  electroplating) memiliki batas semantik yang lebih ambigu, sehingga kedua pakar
  lebih sering tidak sepakat.

  \item \textbf{Penengahan \textit{LLM-as-a-judge}.}

  Mekanisme \textit{LLM-as-a-judge} (Gemini~3.1 Pro Preview) menengahi dua jenis
  ketidaksepakatan intra-mapel pada Fase~1 yang tidak terselesaikan oleh
  kesepakatan langsung antarpakar: (a) perbedaan rating atas \textit{triple} yang
  sama, dan (b) usulan \textit{missing triple} yang hanya diajukan satu pakar.
  Volume dan hasil penengahannya dirangkum pada Tabel~\ref{tab:llm-judge}.

  Untuk perbedaan rating, AI Judge menengahi total 123 \textit{triple} (Biologi
  10, Fisika 55, Kimia 58). Pada 117 dari 123 kasus (95{,}1\%) ia memilih salah
  satu label yang telah diberikan oleh kedua pakar, dan hanya pada 6 kasus
  menetapkan label ketiga yang berbeda. Pola ini menunjukkan bahwa AI Judge
  dominan \emph{menengahi} di antara posisi kedua pakar, bukan melakukan
  penilaian ulang secara independen. Untuk usulan \textit{missing triple}, AI Judge
  menyetujui 29 dari 44 usulan (Fisika 10/15; Kimia 19/29; Biologi tanpa usulan)
  sebagai relasi esensial yang layak ditambahkan ke \textit{ground truth}, dan
  menolak 15 sisanya. Jumlah usulan yang disetujui inilah yang muncul sebagai
  kolom \textbf{Hilang} pada Tabel~\ref{tab:validasi-mapel} dan diperhitungkan
  sebagai \textit{false negative} dalam \textit{Recall}. Putusan lengkap
  AI Judge per \textit{triple} diberikan pada \ref{lamp:missing-triples}.

  Volume penengahan selaras dengan tingkat kesepakatan pada Tabel~\ref{tab:ac1}:
  Biologi yang ber-AC1 tertinggi (0{,}931) hanya menyisakan 10 perbedaan rating,
  sedangkan Fisika dan Kimia yang ber-AC1 lebih rendah menyisakan masing-masing
  55 dan 58 perbedaan yang memerlukan penengahan.
\end{enumerate}

\begin{table}[htbp]
\centering
\caption{Volume dan hasil penengahan \textit{LLM-as-a-judge} atas ketidaksepakatan
  intra-mapel (Fase~1). Kolom \textit{missing} disetujui memuat usulan yang
  diterima sebagai relasi esensial (label \texttt{missing}), sedangkan kolom
  ditolak memuat usulan yang diabaikan (label \texttt{abaikan}).}
\label{tab:llm-judge}
\begin{tabular}{@{}l c c c c@{}}
\toprule
\textbf{Mapel} & \textbf{Perbedaan} & \textbf{Usulan} & \textbf{\textit{Missing}} & \textbf{\textit{Missing}} \\
               & \textbf{rating}    & \textbf{\textit{missing}} & \textbf{disetujui} & \textbf{ditolak} \\
\midrule
Biologi     & 10  & 0  & 0  & 0  \\
Fisika      & 55  & 15 & 10 & 5  \\
Kimia       & 58  & 29 & 19 & 10 \\
\midrule
\textbf{Total} & \textbf{123} & \textbf{44} & \textbf{29} & \textbf{15} \\
\bottomrule
\end{tabular}
\end{table}

% ────────────────────────────────────────────────────────────
\section{Pembahasan Kualitatif Hasil KGC}\label{subsec:kualitatif}
% ────────────────────────────────────────────────────────────

Bagian ini membahas secara kualitatif kualitas relasi lintas-buku yang
ditemukan oleh \textit{pipeline Graph Completion}, menganalisis pola kesalahan
model LLM, dan mengidentifikasi faktor-faktor yang memengaruhi akurasi
ekstraksi.

Pengaruh parameter \textit{retrieval} terhadap jumlah relasi lintas-buku yang
ditemukan dievaluasi pada empat konfigurasi berpasangan (Tabel~\ref{tab:param-sweep});
ambang kemiripan (\textit{threshold}) dinaikkan seiring penurunan jumlah
tetangga per \textit{node} (top-$k$). Semakin ketat konfigurasi, semakin sedikit pasangan
kandidat yang lolos: dari 125 relasi pada konfigurasi paling permisif (0{,}70
/ top-$k$ = 20) menyusut menjadi hanya 14 relasi pada konfigurasi paling ketat
(0{,}85 / top-$k$ = 5). Analisis kualitas per konfigurasi disajikan pada
Subbab~\ref{subsubsec:kualitas}.

\begin{figure}[htbp]
  \centering
  \includegraphics[width=\linewidth]{assets/pics/param-sweep.png}
  \caption{\textit{Parameter sweep}: pengaruh \textit{threshold} dan top-$k$ terhadap
    jumlah relasi lintas-buku yang ditemukan.}
  \label{fig:param-sweep}
\end{figure}

\begin{table}[htbp]
\centering
\caption{Jumlah relasi lintas-buku per konfigurasi \textit{parameter sweep}.}
\label{tab:param-sweep}
\begin{tabular}{@{}
  >{\raggedright\arraybackslash}p{0.40\linewidth}
  >{\raggedright\arraybackslash}p{0.54\linewidth}@{}}
\toprule
\textbf{Konfigurasi (threshold / top-$k$)} & \textbf{Jumlah relasi} \\
\midrule
0{,}70 / $k$=20 (permisif --- konfigurasi final) & 125 \\
0{,}75 / $k$=15                                   & 72  \\
0{,}80 / $k$=10                                   & 42  \\
0{,}85 / $k$=5 (paling ketat)                     & 14  \\
\bottomrule
\end{tabular}
\end{table}

\begin{figure}[htbp]
  \centering
  \scalebox{0.72}{%
  \begin{tikzpicture}[
    layer/.style  = {draw=white, line width=1pt, rounded corners=6pt},
    title/.style  = {font=\small\bfseries, text=black, align=center},
    titlew/.style = {font=\small\bfseries, text=white, align=center},
    contoh/.style = {font=\scriptsize\itshape, text=black!85, align=center},
    contohw/.style= {font=\scriptsize\itshape, text=white, align=center},
  ]
    % Boxes cascade downward (north-anchored, horizontally centered): each
    % inner config starts below its parent's top band, so every title +
    % "Contoh" example stays in its own visible band and never overlaps.
    % Outer (τ=0,70) — palest blue
    \node[layer, fill=blue!12, anchor=north, minimum width=13cm,
          minimum height=10.4cm] (b70) at (0,0) {};
    \node[title, anchor=north] at ([yshift=-6pt]b70.north)
      {0{,}70 / $k$=20 \textemdash{} 117 relasi};
    \node[contoh, anchor=north, text width=12cm] at ([yshift=-26pt]b70.north)
      {Contoh: Arus Listrik $\to$ \texttt{APLIKASI\_DARI} $\to$ Elektroforesis};

    % Mid-1 (τ=0,75)
    \node[layer, fill=blue!28, anchor=north, minimum width=10.5cm,
          minimum height=8.2cm] (b75) at (0,-1.7) {};
    \node[title, anchor=north] at ([yshift=-6pt]b75.north)
      {0{,}75 / $k$=15 \textemdash{} 72 relasi};
    \node[contoh, anchor=north, text width=9.5cm] at ([yshift=-26pt]b75.north)
      {Contoh: Defleksi Gerak Muatan $\to$ \texttt{PRASYARAT\_UNTUK} $\to$ Elektroforesis};

    % Mid-2 (τ=0,80)
    \node[layer, fill=blue!55, anchor=north, minimum width=8cm,
          minimum height=6cm] (b80) at (0,-3.4) {};
    \node[titlew, anchor=north] at ([yshift=-6pt]b80.north)
      {0{,}80 / $k$=10 \textemdash{} 42 relasi};
    \node[contohw, anchor=north, text width=7cm] at ([yshift=-26pt]b80.north)
      {Contoh: Potensial Elektrode $\to$ \texttt{APLIKASI\_DARI} $\to$ Potensial Listrik};

    % Inner (τ=0,85) — darkest
    \node[layer, fill=blue!80, anchor=north, minimum width=5.5cm,
          minimum height=3.8cm] (b85) at (0,-5.1) {};
    \node[titlew, anchor=north] at ([yshift=-6pt]b85.north)
      {0{,}85 / $k$=5 \textemdash{} 14 relasi};
    \node[contohw, anchor=north, text width=4.6cm] at ([yshift=-26pt]b85.north)
      {Contoh: DNA $\leftrightarrow$ \texttt{SAMA\_DENGAN} $\leftrightarrow$ DNA};
  \end{tikzpicture}}

  \caption{Struktur bersarang (\textit{nested}) hasil \textit{parameter sweep}
    relasi lintas-buku antar-konfigurasi.}
  \label{fig:nested-sweep}
\end{figure}

Konfigurasi membentuk himpunan bersarang: setiap konfigurasi yang lebih
ketat merupakan himpunan bagian dari yang lebih permisif. Pengetatan
parameter dengan demikian berfungsi sebagai kompromi \textit{recall} ---
menukar cakupan relasi (lebih banyak interkoneksi ditemukan) dengan
selektivitas (hanya pasangan berkemiripan tertinggi yang dipertahankan).
Sebagaimana akan dibahas pada Subbab~\ref{subsubsec:kualitas}, pengetatan ini menurunkan
kuantitas relasi tetapi \textbf{tidak} memperbaiki ketepatan arah relasi.

% ──────────────────────────────────────────
\subsection{Kualitas Relasi Lintas-Buku yang Ditemukan}\label{subsubsec:kualitas}
% ──────────────────────────────────────────

Pada konfigurasi final (threshold = 0{,}70; top-$k$ = 20; \textit{scope}
lintas-mapel), \textit{pipeline} menemukan 125 relasi lintas-buku. Distribusi
tipe relasi (Gambar~\ref{fig:relasi-distribusi}) didominasi
\texttt{PRASYARAT\_\allowbreak UNTUK} (57 \textit{edge}, 45{,}6\%), diikuti
\texttt{MEMPERDALAM} (31 \textit{edge}, 24{,}8\%),
\texttt{BERKAITAN\_\allowbreak DENGAN} (21 \textit{edge}, 16{,}8\%),
\texttt{APLIKASI\_\allowbreak DARI} (15 \textit{edge}, 12{,}0\%), dan
\texttt{SAMA\_\allowbreak DENGAN} (1 \textit{edge}, 0{,}8\%).
Pada bagian hasil, prefiks \texttt{LINTAS\_BUKU\_} dihilangkan dari nama tipe
relasi untuk menjaga keterbacaan tabel dan gambar; label yang disimpan pada
graf tetap menggunakan bentuk lengkap \texttt{LINTAS\_BUKU\_*} sesuai ontologi.
Dominasi \texttt{PRASYARAT\_\allowbreak UNTUK} mengindikasikan bahwa jembatan lintas-disiplin
pada kurikulum sains kelas XII lebih banyak bersifat prasyarat pengetahuan
daripada aplikasi atau analogi. Dari 125 relasi, 8 \textit{edge} memiliki \textit{confidence}
$< 0{,}5$ (seluruhnya \texttt{BERKAITAN\_\allowbreak DENGAN}) dikeluarkan dari survei validasi
utama, menyisakan 117 relasi yang dinilai pakar.

\begin{figure}[htbp]
  \centering
  \begin{minipage}[c]{0.50\linewidth}
    \centering
    \begin{tikzpicture}
      \begin{axis}[
        xbar,
        width=\linewidth,
        height=5.2cm,
        symbolic y coords={SAMA, BERKAITAN, APLIKASI, MEMPERDALAM, PRASYARAT},
        ytick=data,
        nodes near coords,
        nodes near coords style={font=\scriptsize, /pgf/number format/precision=0, /pgf/number format/fixed},
        every node near coord/.append style={anchor=west, xshift=2pt},
        xmin=0, xmax=70,
        /pgf/bar width=10pt,
        xlabel={Jumlah \textit{edge}},
        xlabel style={font=\scriptsize},
        ticklabel style={font=\scriptsize},
        title={\small Distribusi tipe (total 125)},
        enlarge y limits=0.18,
      ]
      \addplot[fill=blue!45, draw=blue!70] coordinates {
        (1,SAMA)
        (21,BERKAITAN)
        (15,APLIKASI)
        (31,MEMPERDALAM)
        (57,PRASYARAT)
      };
      \end{axis}
    \end{tikzpicture}
  \end{minipage}\hfill
  \begin{minipage}[c]{0.46\linewidth}
    \centering\tiny
    \renewcommand{\arraystretch}{1.3}
    \begin{tabular}{@{}p{1.95cm}@{\,}p{1.55cm}@{\,}p{1.65cm}@{\,}c@{}}
      \toprule
      \textbf{Tipe} & \textbf{Sumber} & \textbf{Target} & \textbf{Conf.} \\
      \midrule
      PRASYARAT\_UNTUK   & DNA (Kim)            & Gen (Bio)            & 1{,}0 \\
      MEMPERDALAM        & DNA (Kim)            & Materi Genetik (Bio) & 0{,}9 \\
      APLIKASI\_DARI     & DNA (Bio)            & Nukleotida (Kim)     & 0{,}9 \\
      BERKAITAN\_DENGAN  & alat ukur listrik (Fis) & jembatan garam (Kim) & 0{,}7 \\
      SAMA\_DENGAN       & DNA (Kim)            & DNA (Bio)            & 1{,}0 \\
      \bottomrule
    \end{tabular}
  \end{minipage}
  \caption{Distribusi tipe relasi lintas-buku hasil klasifikasi LLM beserta contoh triple per tipe.}
  \label{fig:relasi-distribusi}
\end{figure}

Dari 125 \textit{edge}, 20 (16\%) memiliki skor \textit{confidence} = 1{,}0, menunjukkan
keyakinan tinggi LLM terhadap relasi tersebut. Contoh relasi dengan
\textit{confidence} tertinggi meliputi:

\begin{itemize}
  \item DNA (Asam 2-Deoksiribonukleat) [Kimia] $\to$ Gen [Biologi]:
    \texttt{PRASYARAT\_UNTUK} (conf = 1{,}0). Pemahaman struktur kimia DNA merupakan
    fondasi untuk memahami fungsi gen secara biologis.
  \item Potensial Elektrode [Kimia] $\to$ Potensial Listrik [Fisika]:
    \texttt{APLIKASI\_DARI} (conf = 1{,}0). Konsep potensial elektrode dalam
    elektrokimia merupakan aplikasi langsung dari konsep potensial listrik
    dalam fisika.
  \item Enzim [Biologi] $\to$ Protein [Kimia]: \texttt{APLIKASI\_DARI} (conf = 1{,}0).
    Enzim merupakan manifestasi fungsional dari protein dalam konteks biologis.
  \item Fosforilasi Oksidatif [Biologi] $\to$ Oksidasi [Kimia]:
    \texttt{PRASYARAT\_UNTUK} (conf = 0{,}9). Pemahaman reaksi oksidasi secara kimia
    diperlukan sebelum mempelajari fosforilasi oksidatif dalam metabolisme.
\end{itemize}

Satu-satunya relasi bertipe \texttt{SAMA\_DENGAN} yang ditemukan adalah DNA (Asam
2-Deoksiribonukleat) [Kimia] $\leftrightarrow$ DNA (Deoxyribonucleic Acid)
[Biologi] dengan \textit{confidence} 1{,}0. Temuan ini mengonfirmasi bahwa
\textit{pipeline} mampu mendeteksi entitas identik yang diajarkan di dua disiplin
berbeda dengan terminologi yang sedikit berbeda.

\subsubsection{Validasi Pakar atas Relasi Lintas-Buku}

Seluruh 117 relasi kandidat pada konfigurasi final ($\tau$=0{,}70, $k$=20)
divalidasi pakar melalui instrumen survei Fase~2 yang menilai validitas,
ketepatan tipe, dan ketepatan arah tiap relasi secara \textit{score-blind}
melalui diskusi panel; instrumen dan perlakuan metriknya didefinisikan pada
Subbab~\ref{subsec:metrik-fase2}. Hasil per-aspek menunjukkan
\textbf{validitas 99{,}1\%} (116/117), \textbf{ketepatan tipe 98{,}3\%}
(115/117; dua relasi salah tipe: LB008 dan LB074), namun \textbf{ketepatan
arah 75{,}9\%} (88/116; 28 \textit{edge} terbalik dan 1 \textit{edge} dinilai
NA, yakni LB085 yang arahnya tidak dapat ditentukan panel karena keterbatasan
konteks).

Temuan utama validasi ini adalah bahwa \textbf{arah relasi merupakan
sumber kesalahan dominan}, bukan keberadaan atau tipe relasi. Sistem
sangat andal dalam \textit{menemukan} keterkaitan lintas-buku yang nyata
(99{,}1\%) dan dalam \textit{melabeli} tipenya (98{,}3\%), tetapi salah
menentukan arah pada hampir satu dari empat relasi. Kesalahan arah ini
bersifat sistematis: 26 dari 28 relasi terbalik berasal dari pasangan
Biologi\,$\leftrightarrow$\,Kimia: LLM cenderung menempatkan konsep
Biologi sebagai sumber, padahal arah yang benar umumnya
Kimia$\to$Biologi --- konsep kimia berperan sebagai prasyarat/fondasi
yang \textit{diperdalam} atau \textit{dibutuhkan} oleh proses biologis.
Tipe \texttt{MEMPERDALAM} paling rentan (ketepatan arah 60\%: 18 Benar, 12
Terbalik), disusul \texttt{APLIKASI\_DARI} (73\%) dan
\texttt{PRASYARAT\_UNTUK} (79\%),
sementara tipe simetris \texttt{BERKAITAN\_DENGAN} dan \texttt{SAMA\_DENGAN} mencapai
100\% karena arah tidak relevan secara semantik.

Karena konfigurasi yang lebih ketat merupakan himpunan bagian dari
konfigurasi paling permisif (Gambar~\ref{fig:nested-sweep}), \textit{edge}
yang lolos pada $\tau$=0{,}75; 0{,}80; dan 0{,}85 sudah ikut
divalidasi pakar pada survei Fase~2. Hasil agregat per
konfigurasi disajikan pada Tabel~\ref{tab:validasi-sweep}.
Validitas dan ketepatan tipe meningkat dari 99{,}1\% / 98{,}3\%
pada $\tau$=0{,}70 menjadi 100\% pada ketiga konfigurasi yang lebih
ketat, mengonfirmasi bahwa pengetatan ambang menyaring sisa
\textit{noise} validitas/tipe yang ada pada konfigurasi permisif (yakni
LB005, LB008, dan LB074). Namun, ketepatan arah relasi tetap berada di
rentang 71--86\% lintas-konfigurasi tanpa kecenderungan monotonik
terhadap pengetatan parameter, menegaskan bahwa kesalahan arah bukan
persoalan ambang. Relasi terbalik tetap merupakan pasangan
berkemiripan tinggi yang valid secara semantik, sehingga
penanganannya memerlukan perbaikan pada \textit{prompt} klasifikasi
atau normalisasi arah pasca-proses, bukan penyetelan
\textit{threshold}.

{\setlength{\tabcolsep}{4pt}
\begin{table}[htbp]
\centering
\caption{Hasil validasi pakar atas relasi lintas-buku per konfigurasi parameter.
  Kolom \textbf{Edge dihasilkan} adalah jumlah \textit{edge} keluaran
  \textit{sweep}; kolom $n$ adalah jumlah \textit{edge} yang
  benar-benar dinilai pakar pada survei Fase~2.}
\label{tab:validasi-sweep}
\begin{tabular}{@{}
  >{\raggedright\arraybackslash}p{0.20\linewidth}
  >{\centering\arraybackslash}p{0.10\linewidth}
  >{\centering\arraybackslash}p{0.07\linewidth}
  >{\centering\arraybackslash}p{0.16\linewidth}
  >{\centering\arraybackslash}p{0.16\linewidth}
  >{\centering\arraybackslash}p{0.16\linewidth}@{}}
\toprule
\textbf{Konfigurasi ($\tau$ / $k$)} & \textbf{Edge dihasilkan} & \textbf{$n$} & \textbf{Validitas} & \textbf{Ketepatan tipe} & \textbf{Ketepatan arah} \\
\midrule
0{,}70 / $k$=20 (paling permisif) & 125 & 117 & 99{,}1\% (116/117) & 98{,}3\% (115/117) & 75{,}9\% (88/116) \\
0{,}75 / $k$=15                   & 72  & 72  & 100{,}0\% (72/72)  & 100{,}0\% (72/72)  & 75{,}0\% (54/72)  \\
0{,}80 / $k$=10                   & 42  & 42  & 100{,}0\% (42/42)  & 100{,}0\% (42/42)  & 71{,}4\% (30/42)  \\
0{,}85 / $k$=5 (paling ketat)     & 14  & 14  & 100{,}0\% (14/14)  & 100{,}0\% (14/14)  & 85{,}7\% (12/14)  \\
\bottomrule
\end{tabular}
\end{table}
}

\noindent\textit{Catatan.} Pada konfigurasi $\tau$=0{,}70/$k$=20, 8
\textit{edge} dengan \textit{confidence} LLM $< 0{,}5$ (seluruhnya
\texttt{BERKAITAN\_DENGAN}) dikecualikan dari survei utama
(lihat Subbab~\ref{subsec:kualitatif}), sehingga $n=117 < 125$.
Konfigurasi yang lebih ketat ($\tau\!\ge\!0{,}75$) merupakan
himpunan bagian dari subset 117 yang sama, sehingga
\textbf{Edge dihasilkan} = $n$ pada baris kedua dan seterusnya.

\subsubsection{Pemilihan Konfigurasi Final}\label{subsubsec:konfigurasi-final}

Berdasarkan Tabel~\ref{tab:validasi-sweep}, \textbf{konfigurasi final
yang dipilih untuk \textit{Knowledge Graph} adalah
$\tau = 0{,}70$ dengan top-$k = 20$}, menghasilkan 125 relasi lintas-buku
(117 di antaranya divalidasi pakar) dengan validitas 99{,}1\%, ketepatan
tipe 98{,}3\%, dan ketepatan arah 75{,}9\%. Pilihan ini dapat dijustifikasi
melalui tiga pertimbangan yang ditopang oleh data validasi:

\begin{enumerate}
  \item \textbf{Tujuan KGC adalah memaksimalkan cakupan jembatan
    lintas-disiplin}, bukan memaksimalkan presisi relasi. Pengetatan
    \textit{threshold} dari 0{,}70 ke 0{,}75 menaikkan validitas hanya
    0{,}9 poin (99{,}1\% $\to$ 100\%) dan ketepatan tipe 1{,}7 poin
    (98{,}3\% $\to$ 100\%), tetapi sekaligus menghilangkan 45
    \textit{edge} (35\% jembatan lintas-disiplin) yang sebagian besar
    juga valid menurut pakar. Tradeoff ini tidak menguntungkan dalam
    konteks penelitian yang menargetkan keterhubungan antar-mapel
    sebagai luaran utama.

  \item \textbf{\textit{Noise} sisa pada $\tau=0{,}70$ bersifat terbatas
    dan dapat di-\textit{filter} pasca-survei.} Total 3 \textit{edge}
    bermasalah pada konfigurasi paling permisif --- LB005
    (\textit{relasi valid?} = \textit{Tidak}), serta LB008 dan LB074
    (tipe relasi salah) --- teridentifikasi melalui validasi pakar Fase~2
    dan dapat dihapus secara terarah bila diinginkan, tanpa mempengaruhi
    114 \textit{edge} tersurvei lainnya.

  \item \textbf{Pengetatan \textit{threshold} tidak menyelesaikan sumber
    kesalahan dominan.} Sebagaimana telah ditunjukkan sebelumnya, ketepatan
    arah tetap stagnan lintas-konfigurasi (Tabel~\ref{tab:validasi-sweep}
    dan Tabel~\ref{tab:conf-kalibrasi}), sehingga konfigurasi yang lebih
    ketat kehilangan cakupan tanpa imbal balik perbaikan arah.
\end{enumerate}

Konfigurasi alternatif ($\tau = 0{,}75 / k=15$ dengan 72 relasi;
$\tau = 0{,}80 / k=10$ dengan 42 relasi; $\tau = 0{,}85 / k=5$ dengan
14 relasi) tetap relevan sebagai \textit{view} alternatif atas graf
yang sama: peneliti turunan yang memprioritaskan presisi atas cakupan
dapat menyaring graf final menggunakan ambang \textit{cosine}
yang lebih ketat tanpa perlu menjalankan ulang \textit{pipeline}
ekstraksi.

% ──────────────────────────────────────────
\subsection{Analisis \textit{False Positive} dan Kalibrasi \textit{Confidence}}\label{subsubsec:false-positive}
% ──────────────────────────────────────────

Sebanyak 9 \textit{edge} memiliki \textit{confidence} LLM $< 0{,}7$,
seluruhnya bertipe \texttt{BERKAITAN\_DENGAN}; 8 di antaranya bahkan
$< 0{,}5$ dan telah dikecualikan dari survei utama
(Subbab~\ref{subsec:kualitatif}), sehingga hanya 1 \textit{edge} dengan
$0{,}5\!\le\!\text{conf}\!<\!0{,}7$ yang masuk ke validasi pakar. Daftar
berikut menampilkan beberapa contoh berkeyakinan rendah hingga sedang
($\leq 0{,}7$) yang patut dicermati, mencakup baik \textit{edge} yang
dikecualikan maupun yang divalidasi pakar, serta satu \textit{edge}
\texttt{MEMPERDALAM} pada batas $\text{conf}\!=\!0{,}7$ sebagai pembanding:

\begin{itemize}
  \item Gerbang Logika Dasar [Fisika] $\leftrightarrow$ Notasi Sel Standar [Kimia]:
    \texttt{BERKAITAN\_DENGAN} (conf = 0{,}2; \textbf{dikecualikan dari survei}).
    Relasi ini merupakan contoh \textit{false positive} --- kemiripan
    \textit{embedding} terjadi akibat kedua konsep sama-sama melibatkan
    notasi simbolik, namun tidak ada hubungan kurikuler bermakna.
  \item Klasifikasi Semikonduktor [Fisika] $\leftrightarrow$ Polimer Anorganik [Kimia]:
    \texttt{BERKAITAN\_DENGAN} (conf = 0{,}6; \textbf{divalidasi pakar}).
    Kemiripan kemungkinan didorong oleh kesamaan konteks material/bahan,
    namun hubungan substantifnya lemah.
  \item Efek Fotolistrik [Fisika] $\to$ Fotosintesis [Biologi]:
    \texttt{MEMPERDALAM} (conf = 0{,}7; \textbf{divalidasi pakar}; di luar 9
    \textit{edge} \texttt{BERKAITAN\_DENGAN} di atas, ditampilkan sebagai
    pembanding tipe non-\textit{fallback}). Keduanya melibatkan interaksi
    foton dengan materi, tetapi mekanisme dan konteksnya berbeda jauh.
\end{itemize}

Temuan ini mengonfirmasi peran penting label \textit{fallback} \texttt{BERKAITAN\_DENGAN}
sebagai penanda relasi yang lemah atau ambigu. Tipe ini memiliki keyakinan jauh
lebih rendah daripada tipe lain. Rata-rata \textit{confidence} 0{,}50 untuk 21
\textit{edge} \texttt{BERKAITAN\_DENGAN} yang dibandingkan dengan 0{,}85--0{,}94 pada tipe relasi
spesifik menunjukkan bahwa LLM memilih label ini ketika tidak yakin terhadap
arah atau tipe relasi yang lebih spesifik.

Untuk menguji apakah \textit{confidence} LLM merupakan sinyal yang
dapat diandalkan, dilakukan analisis kalibrasi \textit{post hoc}
dengan mem-\textit{bucket} 117 \textit{edge} tersurvei berdasarkan
nilai \textit{confidence} dan menghitung tingkat persetujuan pakar
per \textit{bucket}. Hasilnya disajikan pada
Tabel~\ref{tab:conf-kalibrasi}.

\begin{table}[htbp]
\centering
\caption{Kalibrasi \textit{post hoc} \textit{confidence} LLM terhadap
  penilaian pakar (Fase~2, $n = 117$). Penyebut \textit{ketepatan arah}
  mengecualikan satu \textit{edge} ber-status NA (LB085).}
\label{tab:conf-kalibrasi}
\begin{tabular}{@{}
  >{\raggedright\arraybackslash}p{0.18\linewidth}
  >{\centering\arraybackslash}p{0.07\linewidth}
  >{\centering\arraybackslash}p{0.20\linewidth}
  >{\centering\arraybackslash}p{0.22\linewidth}
  >{\centering\arraybackslash}p{0.25\linewidth}@{}}
\toprule
\textbf{Bucket \textit{conf}} & \textbf{$n$} & \textbf{Validitas} & \textbf{Ketepatan tipe} & \textbf{Ketepatan arah} \\
\midrule
$< 0{,}5$        &  0 & --- & --- & --- \\
$[0{,}5;\;0{,}7)$ &  1 & 1/1 = 100\% & 1/1 = 100\% & 1/1 = 100\% \\
$[0{,}7;\;0{,}85)$ & 27 & 26/27 = 96{,}3\% & 27/27 = 100\% & 20/26 = 76{,}9\% \\
$[0{,}85;\;1{,}0)$ & 69 & 69/69 = 100\% & 67/69 = 97{,}1\% & 52/69 = 75{,}4\% \\
$= 1{,}0$        & 20 & 20/20 = 100\% & 20/20 = 100\% & 15/20 = 75{,}0\% \\
\bottomrule
\end{tabular}
\end{table}

Pola kalibrasi pada Tabel~\ref{tab:conf-kalibrasi} memberikan tiga
implikasi metodologis. \textbf{Pertama}, \textit{confidence}
memiliki sinyal yang cukup baik untuk \textit{keberadaan} relasi:
validitas pakar mencapai 100\% pada seluruh \textit{bucket}
$\geq 0{,}85$ dan tetap di atas 96\% pada \textit{bucket}
$[0{,}7;\;0{,}85)$, sehingga ambang seleksi
(\textsc{berkaitan\_dengan} $\geq 0{,}85$; tipe lain $\geq 0{,}70$;
lihat \ref{lamp:prompt-completion}) terbukti efektif
menyaring \textit{false positive}. \textbf{Kedua}, sinyal terhadap
\textit{ketepatan tipe} hanya marginal: kenaikan dari 97{,}1\%
pada \textit{bucket} $[0{,}85;\;1{,}0)$ ke 100\% pada
\textit{bucket} $\text{conf}=1{,}0$ kecil dan terjadi pada $n$ kecil.
\textbf{Ketiga}, dan terpenting, \textit{confidence} \textbf{tidak}
memberikan sinyal terhadap \textit{ketepatan arah}: ketepatan
arah stagnan pada 75--77\% di seluruh \textit{bucket}
\textit{confidence} ($n > 1$), termasuk pada \textit{bucket}
$\text{conf} = 1{,}0$ yang masih meleset arah pada 5 dari 20
\textit{edge}. Temuan ini secara empiris memperkuat argumen pada
Subbab~\ref{subsubsec:kualitas} bahwa kesalahan arah bukan
persoalan ambang \textit{confidence} --- LLM dapat sepenuhnya yakin
sambil tetap salah arah, sehingga penanganannya memerlukan
perbaikan pada \textit{prompt} klasifikasi atau normalisasi arah
pasca-proses, bukan penyetelan \textit{threshold}. Pola ini juga
sejalan dengan literatur yang menunjukkan keterbatasan kalibrasi
\textit{self-reported confidence} pada model bahasa
\parencite{kadavath_language_2022}.

% ──────────────────────────────────────────
\subsection{Analisis Pola Kesalahan Ekstraksi (Error Analysis)}\label{subsubsec:error-analysis}
% ──────────────────────────────────────────

Analisis \textit{feedback} pakar mengungkapkan empat pola kesalahan sistematis
pada \textit{pipeline} ekstraksi LLM.

\subsubsection{Kesalahan faktual agen biologis.} Pada Kimia, \textit{triple} mengenai produksi
etanol menyatakan ``bakteri memfermentasi karbohidrat'' padahal seharusnya
``ragi (\textit{Saccharomyces cerevisiae})''. Kesalahan ini termasuk dalam kategori
\textit{hallucination} faktual \parencite{ji_survey_2023, min_factscore_2023}, sebagaimana
diidentifikasi pada Subbab~\ref{subsec:kg-ontologi}.

\subsubsection{Pencampuran konsep teknis.} \textit{Triple} mengenai pemurnian tembaga menyatakan
bahwa \textit{electroplating} digunakan untuk pemurnian tembaga, padahal proses yang
benar adalah \textit{elektrorefining}. Kedua konsep memang melibatkan elektrolisis,
namun tujuan dan mekanismenya berbeda.

\subsubsection{Definisi yang terlalu luas.} Expert 2 (Fisika) mencatat sejumlah definisi yang
menimbulkan bias, misalnya definisi konstanta Coulomb yang ``terlalu luas'' tanpa
menyebutkan jenis konstantanya. Umpan balik umum dari pakar ini menekankan:
``Penjelasannya dipastikan jangan ada bias\ldots Lebih ke penulisan, notasi, dan
definisi nya perlu dicek dan dirapihkan kembali.''

\subsubsection{\textit{Under-extraction} rumus matematis.} Expert 1 (Fisika) mengusulkan
15 \textit{missing triples}; 10 di antaranya (67\%) merupakan rumus matematis
yang tidak berhasil diekstrak oleh LLM. Ini menunjukkan bahwa LLM cenderung
mengekstrak konsep deskriptif namun melewatkan relasi kuantitatif berbasis
persamaan matematika. Contoh: rumus efisiensi transformator ($\eta$), nilai
efektif arus AC ($I_\text{ef}$), panjang gelombang minimum sinar-X
($\lambda_\text{min}$), dan rumus peluruhan radioaktif ($N_t$). Secara umum,
sebagian besar komentar Expert 1 (Fisika) bersifat konfirmasi bahwa keluaran
sistem sudah benar, sedangkan koreksi terfokus pada notasi dan rumus.

% ────────────────────────────────────────────────────────────
\section{Analisis Komparatif Antar-Domain Sains}\label{subsec:komparatif}
% ────────────────────────────────────────────────────────────

Bagian ini membandingkan ketiga domain sains dari dua sisi: \textbf{pola
interkoneksi} antar-mata pelajaran (berapa banyak jembatan lintas-buku
yang terbentuk antar pasangan disiplin) dan \textbf{kualitas ekstraksi}
per domain.

\begin{figure}[htbp]
  \centering
  \begin{tikzpicture}[
    mapel/.style={circle, draw=black!65, thick, minimum size=1.8cm,
      align=center, font=\small\bfseries},
    wlab/.style={midway, fill=white, inner sep=2pt, font=\footnotesize}
  ]
    \node[mapel, fill=green!12]  (bio) at (0,0)   {Biologi};
    \node[mapel, fill=blue!12]   (fis) at (6.4,0) {Fisika};
    \node[mapel, fill=orange!20] (kim) at (3.2,3.4) {Kimia};

    \draw[line width=5pt,   draw=black!45] (bio) -- (kim) node[wlab]{91 (78\%)};
    \draw[line width=2pt,   draw=black!45] (kim) -- (fis) node[wlab]{21 (18\%)};
    \draw[line width=0.6pt, draw=black!45] (bio) -- (fis) node[wlab]{5 (4\%)};
  \end{tikzpicture}
  \caption{\textit{Inter-domain bridge graph}: jumlah relasi lintas-buku antar
    pasangan mata pelajaran (117 relasi disurvei; tebal \textit{edge} $\propto$ jumlah).}
  \label{fig:bridge-graph}
\end{figure}

Distribusi jembatan lintas-disiplin sangat tidak merata. Pasangan
Biologi--Kimia mendominasi (91 relasi, 78\%), diikuti Fisika--Kimia (21
relasi, 18\%), sedangkan Biologi--Fisika nyaris tidak terhubung (5
relasi, 4\%). Pola ini menempatkan Kimia sebagai poros (\textit{hub}) yang
menjembatani kedua disiplin lain. Temuan ini konsisten secara pedagogis
\parencite{dyachenko_integration_2024}: tumpang-tindih molekuler dan
biokimiawi (DNA, protein, enzim, reaksi redoks) membuat Biologi dan Kimia
berbagi banyak konsep dasar --- analisis ko-sitasi antar-disiplin
menempatkan biokimia sebagai bidang paling terkoneksi dan kimia pada posisi
sentral di antara sains \parencite{balaban_central_2006} --- sementara
keterkaitan Fisika sebagian besar terbatas pada elektrokimia (arus dan
potensial listrik $\leftrightarrow$ sel elektrokimia). Sifat Fisika yang lebih
abstrak-matematis menjadikannya domain paling terisolasi dalam jaringan
konsep kurikulum kelas XII.

Dari sisi kualitas ekstraksi (validasi Fase 1), terdapat perbedaan
antar-domain:

\begin{table}[htbp]
\centering
\caption{Presisi ekstraksi per domain sains (Fase 1 validasi pakar).}
\label{tab:komparatif-domain}
\begin{tabular}{@{}
  >{\raggedright\arraybackslash}p{0.15\linewidth}
  >{\raggedright\arraybackslash}p{0.50\linewidth}
  >{\raggedright\arraybackslash}p{0.30\linewidth}@{}}
\toprule
\textbf{Domain} & \textbf{Presisi ekstraksi (Fase 1)} & \textbf{Catatan} \\
\midrule
Biologi & $\sim$0{,}95--0{,}97 & paling konsisten; konsep deskriptif-naratif \\
Fisika  & $\sim$0{,}87--1{,}00 & rentang lebar; lemah pada relasi kuantitatif/rumus \\
Kimia   & $\sim$0{,}82--0{,}88 & kesalahan terbanyak pada terminologi proses \\
\bottomrule
\end{tabular}
\end{table}

Perbedaan ini selaras dengan analisis kesalahan pada Subbab~\ref{subsubsec:error-analysis}:
kelemahan utama pada Fisika adalah \textbf{\textit{under-extraction} rumus} (67\%
\textit{missing triples} Fisika berupa formula), sedangkan Kimia rentan pada
\textbf{pencampuran konsep teknis}. Biologi---yang materinya lebih
bersifat deskriptif---paling sesuai dengan kekuatan LLM dalam mengekstrak
struktur konseptual.

\ifSubfilesClassLoaded{\printbibliography[title={Daftar Pustaka}]}{}

\end{document}
